\section{Обнаружение столкновений}

Обнаружение столкновений (collision detection) --- задача определения того, пересекаются ли геометрические 
объекты, и вычисления данных контакта: точки приложения, нормали и глубины проникновения. Решение этой задачи 
является необходимым условием для корректной обработки контактных взаимодействий в физическом симуляторе.

В практических симуляторах задача решается в два этапа~\cite{ericson2004real}: \emph{широкая фаза} 
(broadphase) быстро отсеивает пары тел, между которыми столкновение невозможно, а \emph{узкая фаза} 
(narrowphase) точно вычисляет контакт для оставшихся пар.

\subsection{Широкая фаза: ограничивающие объёмы}

\subsubsection{Axis-Aligned Bounding Box (AABB)}

AABB --- прямоугольный параллелепипед с рёбрами, параллельными осям координат. Для каждого тела хранятся 
минимальная и максимальная точки по каждой оси: $(\mathbf{p}_{\min},\, \mathbf{p}_{\max})$.

Проверка перекрытия двух AABB выполняется за $O(1)$ с помощью теоремы о разделяющей оси:

\begin{equation}
  \text{overlap}(A, B) \iff
  A_{\max}^x \geq B_{\min}^x \;\wedge\; B_{\max}^x \geq A_{\min}^x \;\wedge\;
  A_{\max}^y \geq B_{\min}^y \;\wedge\; \ldots
  \label{eq:aabb-test}
\end{equation}

При движении тел AABB пересчитывается на каждом шаге. Для ускорения пересчёта применяют раздутые AABB (fat 
AABB), которые обновляются только при выходе тела за границы раздутого объёма.

\subsubsection{Sweep and Prune (SAP)}

Алгоритм SAP~\cite{cohen1995sahp} сортирует интервалы AABB вдоль каждой из трёх осей. Пара тел является 
кандидатом на столкновение только если их интервалы перекрываются по всем трём осям.

Сложность обновления: $O(n + k)$ при когерентном движении, где $n$ — число тел, $k$ --- число перемещений 
порядка (обычно $O(n)$ для типичных сцен).

В контексте графического процессора SAP трудно распараллеливается из-за зависимостей при сортировке. 
Альтернативой служит \emph{пространственное хеширование} (spatial hashing): каждое тело помещается в ячейки 
равномерной сетки, кандидаты на столкновение --- тела в одних и тех же ячейках. Операции вставки и поиска 
хорошо распараллеливаются на графическом процессоре~\cite{harada2007physics}.

\subsection{Узкая фаза: алгоритм GJK}

Алгоритм Гилберта–Джонсона–Кирти (GJK)~\cite{gilbert1988fast} находит ближайшие точки двух выпуклых множеств, 
используя понятие \emph{опорной функции} (support function):

\begin{equation}
  \mathbf{s}_A(\mathbf{d}) = \underset{\mathbf{p} \in A}{\arg\max}\; \mathbf{p} \cdot \mathbf{d},
\end{equation}

то есть наиболее удалённую точку множества $A$ в направлении $\mathbf{d}$.

Ключевая идея GJK: тела $A$ и $B$ пересекаются тогда и только тогда, когда их \emph{минковская разность}

\begin{equation}
  A \ominus B = \{ \mathbf{a} - \mathbf{b} \mid \mathbf{a} \in A,\; \mathbf{b} \in B \}
\end{equation}
содержит начало координат.
Опорная функция минковской разности:
\begin{equation}
  \mathbf{s}_{A \ominus B}(\mathbf{d}) = \mathbf{s}_A(\mathbf{d}) - \mathbf{s}_B(-\mathbf{d}).
\end{equation}

Алгоритм строит симплекс (отрезок, треугольник или тетраэдр) в $A \ominus B$ и итеративно приближается к 
ближайшей к началу координат точке. На каждой итерации выполняется одно обращение к опорным функциям, в 
типичных случаях сходимость достигается за 2–4 итерации~\cite{bergen2004game}.

\paragraph{Поддерживаемые примитивы.}
Опорные функции существуют для всех стандартных примитивов:
\begin{itemize}
  \item сфера: $\mathbf{s}(\mathbf{d}) = \mathbf{c} + r\,\hat{\mathbf{d}}$;
  \item капсула: $\mathbf{s}(\mathbf{d}) = \mathbf{p}_{\text{extremal}} + r\,\hat{\mathbf{d}}$, где $\mathbf{p}_{\text{extremal}}$ — конечная точка оси, ближайшая к $\mathbf{d}$;
  \item параллелепипед: $s_i(\mathbf{d}) = c_i + \text{sign}(d_i \cdot \hat{\mathbf{e}}_i)\,h_i$ по каждой полуоси $\hat{\mathbf{e}}_i$ длиной $h_i$;
  \item выпуклый многогранник: $\mathbf{s}(\mathbf{d}) = \underset{\mathbf{v} \in V}{\arg\max}\; \mathbf{v} \cdot \mathbf{d}$ по вершинам $V$.
\end{itemize}

\subsection{Расширенный политоп: алгоритм EPA}

Если GJK обнаружил пересечение (начало координат внутри симплекса), для вычисления нормали контакта и глубины 
проникновения применяется алгоритм расширяющегося политопа (Expanding Polytope Algorithm, 
EPA)~\cite{bergen2004game}.

Идея EPA: начиная с финального симплекса GJK, итеративно расширять выпуклый политоп, добавляя новые точки 
минковской разности, пока не будет найдена ближайшая к началу грань. Нормаль ближайшей грани даёт направление 
проникновения, расстояние до начала — глубину проникновения $d$.

Данные контакта:

\begin{itemize}
  \item $\hat{\mathbf{n}}$ --- нормаль столкновения (направление вывода тела $B$ из $A$);
  \item $d$ --- глубина проникновения ($d > 0$ при пересечении);
  \item $\mathbf{p}_A,\, \mathbf{p}_B$ --- точки контакта на поверхностях тел.
\end{itemize}

\subsection{Устойчивые контакты}

В реальных сценах тела могут находиться в длительном контакте (стопка кубов, опирание манипулятора на стол). 
Если на каждом шаге вычислять новые контактные точки с нуля, возникает «джиттер» --- высокочастотные 
осцилляции из-за численных неточностей при определении активного набора контактов.

Для устранения этого эффекта применяют \emph{устойчивые контакты} (persistent contacts):

\begin{enumerate}
  \item Контактные точки кешируются от шага к шагу.
  \item На новом шаге проверяется, остаётся ли кешированная точка активной (тела всё ещё достаточно близко).
  \item Новые контакты добавляются, недействительные --- удаляются.
  \item Для решателя формируется набор из 4 наиболее значимых контактных точек (по площади опорного многоугольника).
\end{enumerate}

\subsection{Выводы по разделу}

Система обнаружения столкновений строится как двухфазный конвейер: broadphase на основе AABB отсеивает пары 
тел за $O(n \log n)$, narrowphase на основе GJK/EPA точно вычисляет контактные данные за несколько итераций. 
Кеширование устойчивых контактов устраняет джиттер и снижает вычислительную нагрузку. В разделе~2.4 
рассмотрено, как полученные контактные данные используются в решателе ограничений.
